\documentclass[11pt, letterpaper]{article}

% --- Idioma y codificación ---
\usepackage[spanish]{babel}
\usepackage[utf8]{inputenc}
\usepackage[T1]{fontenc}
\usepackage{lmodern}

% --- Formato ---
\usepackage[margin=2cm, landscape]{geometry}
\usepackage{parskip}
\usepackage{booktabs}
\usepackage{array}
\usepackage{multirow}
\usepackage{colortbl}
\usepackage{xcolor}
\usepackage{tikz}
\usetikzlibrary{positioning, arrows.meta, calc, fit, backgrounds, decorations.pathreplacing}
\usepackage{hyperref}
\hypersetup{colorlinks=true, linkcolor=blue!70!black, urlcolor=blue!70!black}
\usepackage{enumitem}
\usepackage{float}

% --- Colores por pista temática ---
\definecolor{cTrack}{HTML}{DCE6F1}   % azul claro — pista C
\definecolor{rTrack}{HTML}{E2EFDA}   % verde claro — pista Rust
\definecolor{lTrack}{HTML}{FFF2CC}   % amarillo claro — pista Linux/Infra
\definecolor{qTrack}{HTML}{F2DCDB}   % rosa claro — pista Calidad/Herramientas
\definecolor{tTrack}{HTML}{E4DFEC}   % lavanda — pista Troncal
\definecolor{dTrack}{HTML}{FDDCC8}   % naranja claro — pista DevOps/Cloud
\definecolor{goTrack}{HTML}{D1F2EB}  % turquesa claro — pista Go
\definecolor{rowhead}{HTML}{4472C4}  % azul oscuro para encabezados

\title{%
  Mapa de Dependencias del Curso\\[4pt]
  \large Linux $\cdot$ C $\cdot$ Rust $\cdot$ Go
}
\author{22 bloques --- B01 a B22}
\date{}

\begin{document}
\maketitle

% ===================================================================
\section{Tabla de flujo del curso}
% ===================================================================

Cada fila representa un \textbf{nivel de profundidad}. Dentro de un
mismo nivel, los bloques son independientes entre sí y pueden cursarse
en cualquier orden (o en paralelo). Para avanzar al nivel siguiente se
necesitan los prerequisitos indicados en la Sección~\ref{sec:deps}.

\bigskip

\renewcommand{\arraystretch}{1.8}
\setlength{\tabcolsep}{3pt}

\begin{table}[H]
\centering
\small
\begin{tabular}{
  >{\centering\arraybackslash}m{1.8cm}|
  >{\centering\arraybackslash}m{2.4cm}|
  >{\centering\arraybackslash}m{2.4cm}|
  >{\centering\arraybackslash}m{2.4cm}|
  >{\centering\arraybackslash}m{2.4cm}|
  >{\centering\arraybackslash}m{2.4cm}|
  >{\centering\arraybackslash}m{2.4cm}|
  >{\centering\arraybackslash}m{2.4cm}}
\toprule
\rowcolor{rowhead}
\textcolor{white}{\textbf{Nivel}}
  & \textcolor{white}{\textbf{Herramientas}}
  & \textcolor{white}{\textbf{Teoría / Datos}}
  & \textcolor{white}{\textbf{Pista C}}
  & \textcolor{white}{\textbf{Pista Rust}}
  & \textcolor{white}{\textbf{Pista Go}}
  & \textcolor{white}{\textbf{Linux / Infra}}
  & \textcolor{white}{\textbf{DevOps / Cloud}} \\
\midrule
%%% Nivel 1 %%%
\cellcolor{tTrack}\textbf{Nivel 1}\newline Fundamentos
  & \cellcolor{tTrack}\textbf{B01}\newline Docker
  & \cellcolor{tTrack}\textbf{B02}\newline Fund.\ Linux
  & \cellcolor{cTrack}\textbf{B03}\newline Fund.\ C
  & \cellcolor{rTrack}\textbf{B05}\newline Fund.\ Rust
  & \cellcolor{goTrack}\textbf{B22}\newline Fund.\ Go
  & ---
  & --- \\
\hline
%%% Nivel 2 %%%
\cellcolor{qTrack}\textbf{Nivel 2}\newline Ramificación
  & \cellcolor{qTrack}\textbf{B04}\newline Sist.\ Compilación
  & \cellcolor{tTrack}\textbf{B15}\newline Estr.\ de Datos
  & \cellcolor{cTrack}\textbf{B06}\newline Prog.\ Sist.\ C
  & \cellcolor{rTrack}\textbf{B07}\newline Prog.\ Sist.\ Rust
  & ---
  & \cellcolor{lTrack}\textbf{B08} Almac.\newline
    \textbf{B09} Redes
  & --- \\
\hline
%%% Nivel 3 %%%
\cellcolor{qTrack}\textbf{Nivel 3}\newline Aplicado
  & \cellcolor{qTrack}\textbf{B16}\newline Patrones
  & \cellcolor{dTrack}\textbf{B19}\newline Bases de Datos
  & ---
  & \cellcolor{rTrack}\textbf{B12} GUI\newline
    \textbf{B13} Multimedia
  & ---
  & \cellcolor{lTrack}\textbf{B10}\newline Servicios Red
  & \cellcolor{dTrack}\textbf{B18}\newline Scripting \\
\hline
%%% Nivel 4 %%%
\cellcolor{qTrack}\textbf{Nivel 4}\newline Avanzado
  & \cellcolor{qTrack}\textbf{B17}\newline Testing+IS
  & \cellcolor{lTrack}\textbf{B11}\newline Seguridad
  & ---
  & \cellcolor{tTrack}\textbf{B14}\newline Interop/WASM
  & ---
  & ---
  & \cellcolor{dTrack}\textbf{B20}\newline Kubernetes \\
\hline
%%% Nivel 5 %%%
\cellcolor{dTrack}\textbf{Nivel 5}\newline Capstone
  & ---
  & ---
  & ---
  & ---
  & ---
  & ---
  & \cellcolor{dTrack}\textbf{B21}\newline Cloud+TF \\
\bottomrule
\end{tabular}
\caption{Flujo del curso por niveles y pistas temáticas.
  Los colores indican:
  \colorbox{tTrack}{troncal},
  \colorbox{cTrack}{pista C},
  \colorbox{rTrack}{pista Rust},
  \colorbox{goTrack}{pista Go},
  \colorbox{lTrack}{Linux/Infra},
  \colorbox{qTrack}{Herramientas},
  \colorbox{dTrack}{DevOps/Cloud}.}
\end{table}

\newpage

% ===================================================================
\section{Prerequisitos detallados}\label{sec:deps}
% ===================================================================

\renewcommand{\arraystretch}{1.5}

\begin{table}[H]
\centering
\begin{tabular}{c l l l l}
\toprule
\textbf{Bloque} & \textbf{Nombre} & \textbf{Requiere} & \textbf{Recomendado} & \textbf{Aporta a} \\
\midrule
\rowcolor{tTrack}
B01 & Docker y Entorno & --- & --- & B02, B20 \\
\rowcolor{tTrack}
B02 & Fundamentos de Linux & B01 & --- & B03, B08, B09, B11, B18, B19 \\
\rowcolor{cTrack}
B03 & Fundamentos de C & B02 & --- & B04, B05, B06, B15, B16 \\
\rowcolor{qTrack}
B04 & Sistemas de Compilación & B03 & --- & B06, B15 \\
\rowcolor{rTrack}
B05 & Fundamentos de Rust & B03 & --- & B07, B12, B13, B15, B16 \\
\rowcolor{cTrack}
B06 & Prog.\ de Sistemas en C & B03, B04 & --- & B08, B10, B11, B14 \\
\rowcolor{rTrack}
B07 & Prog.\ de Sistemas en Rust & B05 & --- & B12, B13, B14 \\
\rowcolor{lTrack}
B08 & Almacenamiento y Sist.\ Arch. & B02 & B06 & --- \\
\rowcolor{lTrack}
B09 & Redes & B02 & --- & B10, B18, B19, B21 \\
\rowcolor{lTrack}
B10 & Servicios de Red & B09 & B06 & --- \\
\rowcolor{lTrack}
B11 & Seguridad, Kernel y Arranque & B02 & B06 & --- \\
\rowcolor{rTrack}
B12 & GUI con Rust (egui) & B05 & B07 & --- \\
\rowcolor{rTrack}
B13 & Multimedia con GStreamer & B05 & B07 & --- \\
\rowcolor{tTrack}
B14 & Interop C/Rust y WASM & B06, B07 & --- & --- \\
\rowcolor{tTrack}
B15 & Estructuras de Datos & B03, B05 & B04 & B16 \\
\rowcolor{qTrack}
B16 & Patrones de Diseño & B03, B05 & B15 & B17 \\
\rowcolor{qTrack}
B17 & Testing e Ing.\ de Software & B03, B05 & B15, B16 & --- \\
\rowcolor{dTrack}
B18 & Scripting y Automatización & B02 & B09 & B21 \\
\rowcolor{dTrack}
B19 & Bases de Datos & B02 & B09, B03, B05 & B21 \\
\rowcolor{dTrack}
B20 & Kubernetes y Observabilidad & B01 & B09, B18, B19 & B21 \\
\rowcolor{dTrack}
B21 & Cloud Computing y Terraform & B09, B18 & B01, B19, B20 & --- \\
\rowcolor{goTrack}
B22 & Fundamentos de Go & B02 & B09 & B20, B21 \\
\bottomrule
\end{tabular}
\caption{Prerequisitos obligatorios y recomendados de cada bloque.}
\end{table}

\newpage

% ===================================================================
\section{Grafo de dependencias}
% ===================================================================

Flechas sólidas = prerequisito obligatorio.
Flechas punteadas = recomendado pero no obligatorio.

\bigskip

\begin{center}
\begin{tikzpicture}[
  node distance=1.2cm and 1.6cm,
  block/.style={
    rectangle, rounded corners=4pt, draw, thick,
    minimum width=2.8cm, minimum height=1cm,
    font=\small\bfseries, align=center,
    text=black
  },
  troncal/.style={block, fill=tTrack},
  cpista/.style={block, fill=cTrack},
  rpista/.style={block, fill=rTrack},
  lpista/.style={block, fill=lTrack},
  qpista/.style={block, fill=qTrack},
  dpista/.style={block, fill=dTrack},
  gopista/.style={block, fill=goTrack},
  req/.style={-{Stealth[length=6pt]}, thick},
  rec/.style={-{Stealth[length=6pt]}, thick, densely dashed, gray},
]

% Constrain bounding box to prevent arrow curves from overflowing the page
\useasboundingbox (-2,-15) rectangle (21,1);

% --- Nivel 1 ---
\node[troncal] (b01) {B01\\Docker};
\node[troncal, right=2.0cm of b01] (b02) {B02\\Linux};
\node[cpista,  right=2.0cm of b02] (b03) {B03\\Fund.\ C};
\node[rpista,  right=2.0cm of b03] (b05) {B05\\Fund.\ Rust};

% Go (Nivel 1, below B02 to the right)
\node[gopista, below=0.6cm of b02] (b22) {B22\\Fund.\ Go};

% --- Nivel 2 ---
\node[qpista,  below=1.8cm of b01] (b04) {B04\\Build};
\node[cpista,  below=1.8cm of b03] (b06) {B06\\Sistemas C};
\node[rpista,  below=1.8cm of b05] (b07) {B07\\Sistemas Rust};
\node[troncal, below=1.8cm of b02] (b15) {B15\\Estr.\ Datos};

% Infra branch (right)
\node[lpista,  right=2.0cm of b05] (b08) {B08\\Almac.};
\node[lpista,  below=0.6cm of b08] (b09) {B09\\Redes};

% --- Nivel 3 ---
\node[lpista,  below=1.8cm of b09] (b10) {B10\\Servicios};
\node[rpista,  below=1.8cm of b07] (b12) {B12\\GUI};
\node[qpista,  below=1.8cm of b04] (b16) {B16\\Patrones};
\node[rpista,  left=0.4cm of b12] (b13) {B13\\Multimedia};
\node[lpista,  below=1.8cm of b15] (b11) {B11\\Seguridad};

% --- Nivel 4 ---
\node[tTrack, block, below=1.8cm of b11] (b14) {B14\\Interop/WASM};
\node[qpista,  below=1.8cm of b16] (b17) {B17\\Testing+IS};

% --- DevOps/SRE track (Nivel 4-6) ---
\node[dpista,  right=3.0cm of b14] (b19) {B19\\BD};
\node[dpista,  right=1.5cm of b19] (b18) {B18\\Scripting};
\node[dpista,  below=1.8cm of b19, xshift=2cm] (b20) {B20\\K8s+Obs.};
\node[dpista,  below=1.8cm of b20] (b21) {B21\\Cloud+TF};

% === Flechas obligatorias ===
\draw[req] (b01) -- (b02);
\draw[req] (b02) -- (b03);
\draw[req] (b03) -- (b05);
\draw[req] (b03) -- (b04);

% B06 needs B03 + B04
\draw[req] (b03) -- (b06);
\draw[req] (b04) -- (b06);

% B07 needs B05
\draw[req] (b05) -- (b07);

% Infra needs B02
\draw[req] (b02.north east) -- (b08.west);
\draw[req] (b02) -- (b09.west);
\draw[req] (b09) -- (b10);

% B15 needs B03 + B05
\draw[req] (b03) -- (b15);
\draw[req] (b05) -- (b15);

% B11 needs B02
\draw[req] (b02) -- (b11);

% B14 needs B06 + B07
\draw[req] (b06) -- (b14);
\draw[req] (b07) -- (b14);

% B16 needs B03 + B05 (via B15 recommended)
\draw[req] (b03.south) to[out=-90,in=160] (b16.north west);
\draw[req] (b05.south) to[out=-90,in=20] (b16.north east);

% B17 needs B03 + B05 (through B16)
\draw[req] (b16) -- (b17);

% B12, B13 need B05
\draw[req] (b05) -- (b12);
\draw[req] (b05) -- (b13);

% B22 needs B02
\draw[req] (b02) -- (b22);

% DevOps obligatorias
\draw[req] (b02.south east) to[out=-20,in=160] (b18.north west);  % B18 ← B02
\draw[req] (b02.south east) to[out=-25,in=150] (b19.north west);  % B19 ← B02
\draw[req] (b01.south) to[out=-60,in=170] (b20.west);             % B20 ← B01
\draw[req] (b09.south east) to[out=-45,in=45] (b21.north east);   % B21 ← B09
\draw[req] (b18.south) to[out=-100,in=30] (b21.north east);       % B21 ← B18

% === Flechas recomendadas ===
\draw[rec] (b06) -- (b08);
\draw[rec] (b06) to[bend right=15] (b10);
\draw[rec] (b06) -- (b11);
\draw[rec] (b07) -- (b12);
\draw[rec] (b07) -- (b13);
\draw[rec] (b15) -- (b16);
\draw[rec] (b15) to[out=-110,in=110] (b17);
\draw[rec] (b04) -- (b15.north west);

% B22 recomendadas
\draw[rec] (b09.west) to[out=180,in=-30] (b22.south east);        % B22 rec B09

% DevOps recomendadas
\draw[rec] (b09.south) to[out=-60,in=50] (b18.north east);        % B18 rec B09
\draw[rec] (b09.south west) to[out=-100,in=50] (b19.north east);  % B19 rec B09
\draw[rec] (b09.south) to[out=-50,in=50] (b20.north east);        % B20 rec B09
\draw[rec] (b18.south) to[out=-110,in=100] (b20.north);           % B20 rec B18
\draw[rec] (b19.south) to[out=-30,in=160] (b20.north west);       % B20 rec B19
\draw[rec] (b01.south west) to[out=-50,in=160] (b21.north west);  % B21 rec B01
\draw[rec] (b19.south) to[out=-80,in=140] (b21.north west);       % B21 rec B19
\draw[rec] (b20) -- (b21);                                         % B21 rec B20

\end{tikzpicture}
\end{center}

\newpage

% ===================================================================
\section{Rutas de estudio sugeridas}
% ===================================================================

\subsection{Ruta troncal (obligatoria para todos)}

\begin{center}
\begin{tikzpicture}[
  node distance=0.6cm,
  blk/.style={rectangle, rounded corners=3pt, draw, thick,
    fill=tTrack, minimum width=1.8cm, minimum height=0.7cm,
    font=\small\bfseries},
  arr/.style={-{Stealth[length=5pt]}, thick}
]
\node[blk] (a) {B01};
\node[blk, right=of a] (b) {B02};
\node[blk, right=of b] (c) {B03};
\node[blk, right=of c] (d) {B05};
\draw[arr] (a)--(b); \draw[arr] (b)--(c); \draw[arr] (c)--(d);
\end{tikzpicture}
\end{center}

Todo lo demás parte de estos cuatro bloques. Son secuenciales y
obligatorios.

\subsection{Ruta SysAdmin / Infraestructura}

\begin{center}
\begin{tikzpicture}[
  node distance=0.6cm,
  blk/.style={rectangle, rounded corners=3pt, draw, thick,
    fill=lTrack, minimum width=1.8cm, minimum height=0.7cm,
    font=\small\bfseries},
  goblk/.style={rectangle, rounded corners=3pt, draw, thick,
    fill=goTrack, minimum width=1.8cm, minimum height=0.7cm,
    font=\small\bfseries},
  trn/.style={rectangle, rounded corners=3pt, draw, thick,
    fill=tTrack, minimum width=1.8cm, minimum height=0.7cm,
    font=\small\bfseries},
  arr/.style={-{Stealth[length=5pt]}, thick}
]
\node[trn] (a) {troncal};
\node[blk, right=of a] (b) {B08};
\node[blk, right=of b] (c) {B09};
\node[blk, right=of c] (d) {B10};
\node[blk, right=of d] (e) {B11};
\node[goblk, right=of e] (f) {B22};
\draw[arr] (a)--(b); \draw[arr] (b)--(c);
\draw[arr] (c)--(d); \draw[arr] (d)--(e); \draw[arr] (e)--(f);
\end{tikzpicture}
\end{center}

Ideal para quien busca certificaciones Linux (LPIC, RHCSA) o
administración de servidores. B06 recomendado antes de B10 y B11.
B22 (Go) puede cursarse tras B09; se lista al final porque complementa
la ruta con el lenguaje dominante del ecosistema SysAdmin (Docker,
Kubernetes, Terraform, Prometheus).

\subsection{Ruta C completa (sistemas + calidad)}

\begin{center}
\begin{tikzpicture}[
  node distance=0.6cm,
  blk/.style={rectangle, rounded corners=3pt, draw, thick,
    fill=cTrack, minimum width=1.8cm, minimum height=0.7cm,
    font=\small\bfseries},
  qblk/.style={rectangle, rounded corners=3pt, draw, thick,
    fill=qTrack, minimum width=1.8cm, minimum height=0.7cm,
    font=\small\bfseries},
  trn/.style={rectangle, rounded corners=3pt, draw, thick,
    fill=tTrack, minimum width=1.8cm, minimum height=0.7cm,
    font=\small\bfseries},
  arr/.style={-{Stealth[length=5pt]}, thick}
]
\node[trn] (a) {troncal};
\node[qblk, right=of a] (b) {B04};
\node[blk, right=of b] (c) {B06};
\node[trn, right=of c] (d) {B15};
\node[qblk, right=of d] (e) {B16};
\node[qblk, right=of e] (f) {B17};
\draw[arr] (a)--(b); \draw[arr] (b)--(c);
\draw[arr] (c)--(d); \draw[arr] (d)--(e); \draw[arr] (e)--(f);
\end{tikzpicture}
\end{center}

Para desarrollo de sistemas, embebidos, o contribución a proyectos
de código abierto en C.

\subsection{Ruta Rust completa (aplicaciones + interop)}

\begin{center}
\begin{tikzpicture}[
  node distance=0.6cm,
  blk/.style={rectangle, rounded corners=3pt, draw, thick,
    fill=rTrack, minimum width=1.8cm, minimum height=0.7cm,
    font=\small\bfseries},
  trn/.style={rectangle, rounded corners=3pt, draw, thick,
    fill=tTrack, minimum width=1.8cm, minimum height=0.7cm,
    font=\small\bfseries},
  arr/.style={-{Stealth[length=5pt]}, thick}
]
\node[trn] (a) {troncal};
\node[blk, right=of a] (b) {B07};
\node[blk, right=of b] (c) {B12};
\node[blk, right=of c] (d) {B13};
\node[trn, right=of d] (e) {B14};
\draw[arr] (a)--(b); \draw[arr] (b)--(c);
\draw[arr] (c)--(d); \draw[arr] (d)--(e);
\end{tikzpicture}
\end{center}

Para desarrollo de aplicaciones en Rust, WebAssembly, o GUI.
B14 (Interop) requiere también B06, así que la ruta C parcial es
necesaria para llegar ahí.

\subsection{Ruta DevOps / SRE}

\begin{center}
\begin{tikzpicture}[
  node distance=0.6cm,
  blk/.style={rectangle, rounded corners=3pt, draw, thick,
    fill=dTrack, minimum width=1.8cm, minimum height=0.7cm,
    font=\small\bfseries},
  lblk/.style={rectangle, rounded corners=3pt, draw, thick,
    fill=lTrack, minimum width=1.8cm, minimum height=0.7cm,
    font=\small\bfseries},
  goblk/.style={rectangle, rounded corners=3pt, draw, thick,
    fill=goTrack, minimum width=1.8cm, minimum height=0.7cm,
    font=\small\bfseries},
  trn/.style={rectangle, rounded corners=3pt, draw, thick,
    fill=tTrack, minimum width=1.8cm, minimum height=0.7cm,
    font=\small\bfseries},
  arr/.style={-{Stealth[length=5pt]}, thick}
]
\node[trn] (a) {troncal};
\node[lblk, right=of a] (b) {B09};
\node[goblk, right=of b] (g) {B22};
\node[blk, right=of g] (c) {B18};
\node[blk, right=of c] (d) {B19};
\node[blk, right=of d] (e) {B20};
\node[blk, right=of e] (f) {B21};
\draw[arr] (a)--(b); \draw[arr] (b)--(g);
\draw[arr] (g)--(c); \draw[arr] (c)--(d);
\draw[arr] (d)--(e); \draw[arr] (e)--(f);
\end{tikzpicture}
\end{center}

Para roles de SysAdmin, DevOps, o SRE. B22 (Go) se recomienda antes de
B18 para poder escribir herramientas propias en Go además de scripts
Bash/Python. Go es especialmente útil antes de B20 (Kubernetes) ya que
K8s, Prometheus, y la mayoría de herramientas cloud-native están escritas
en Go. B01 (Docker) recomendado en paralelo con B09.

\subsection{Ruta WebAssembly (camino más corto a B14)}

\begin{center}
\begin{tikzpicture}[
  node distance=0.6cm,
  blk/.style={rectangle, rounded corners=3pt, draw, thick,
    minimum width=1.8cm, minimum height=0.7cm,
    font=\small\bfseries},
  trn/.style={blk, fill=tTrack},
  cblk/.style={blk, fill=cTrack},
  rblk/.style={blk, fill=rTrack},
  qblk/.style={blk, fill=qTrack},
  ext/.style={rectangle, rounded corners=3pt, draw, thick,
    minimum width=2.4cm, minimum height=0.7cm,
    font=\small\bfseries, fill=gray!15, densely dashed},
  arr/.style={-{Stealth[length=5pt]}, thick},
  par/.style={decorate, decoration={brace, amplitude=5pt, mirror},
    thick, gray}
]
% Main path
\node[trn] (b01) {B01};
\node[trn, right=of b01] (b02) {B02};
\node[cblk, right=of b02] (b03) {B03};
% Fork: C track (top) and Rust track (bottom)
\node[qblk, above right=0.6cm and 0.6cm of b03] (b04) {B04};
\node[rblk, below right=0.6cm and 0.6cm of b03] (b05) {B05};
\node[cblk, right=of b04] (b06) {B06};
\node[rblk, right=of b05] (b07) {B07};
% Merge at B14
\node[trn, below right=0.6cm and 0.6cm of b06] (b14) {B14};
% External goal
\node[ext, right=1.0cm of b14] (wasm) {WASM App};
% External self-study
\node[ext, below=0.8cm of b14] (js) {JS/TS/Canvas};

\draw[arr] (b01)--(b02);
\draw[arr] (b02)--(b03);
\draw[arr] (b03)--(b04);
\draw[arr] (b03)--(b05);
\draw[arr] (b04)--(b06);
\draw[arr] (b05)--(b07);
\draw[arr] (b06)--(b14);
\draw[arr] (b07)--(b14);
\draw[arr] (b14)--(wasm);
\draw[arr, densely dashed, gray] (js)--(wasm);

% Parallel braces
\draw[par] ($(b04.north west)+(-0.1,0.15)$) -- ($(b05.south west)+(-0.1,-0.15)$)
  node[midway, left=6pt, font=\scriptsize, gray, align=center] {en\\paralelo};
\draw[par] ($(b06.north west)+(-0.1,0.15)$) -- ($(b07.south west)+(-0.1,-0.15)$)
  node[midway, left=6pt, font=\scriptsize, gray, align=center] {en\\paralelo};
\end{tikzpicture}
\end{center}

Camino más corto hacia WebAssembly: \textbf{7 bloques}. B04+B05 pueden
cursarse en paralelo, y después B06+B07 también, reduciendo el camino
efectivo a \textbf{5 fases}. JS, TypeScript y Shaders/Canvas se estudian
por cuenta propia en paralelo con cualquier fase del curso.

\subsection{Ruta Go / Tooling SysAdmin}

\begin{center}
\begin{tikzpicture}[
  node distance=0.6cm,
  blk/.style={rectangle, rounded corners=3pt, draw, thick,
    minimum width=1.8cm, minimum height=0.7cm,
    font=\small\bfseries},
  trn/.style={blk, fill=tTrack},
  goblk/.style={blk, fill=goTrack},
  lblk/.style={blk, fill=lTrack},
  dblk/.style={blk, fill=dTrack},
  arr/.style={-{Stealth[length=5pt]}, thick}
]
\node[trn] (b01) {B01};
\node[trn, right=of b01] (b02) {B02};
\node[goblk, right=of b02] (b22) {B22};
\node[lblk, right=of b22] (b09) {B09};
\node[dblk, right=of b09] (b18) {B18};
\node[dblk, right=of b18] (b20) {B20};
\node[dblk, right=of b20] (b21) {B21};
\draw[arr] (b01)--(b02); \draw[arr] (b02)--(b22);
\draw[arr] (b22)--(b09); \draw[arr] (b09)--(b18);
\draw[arr] (b18)--(b20); \draw[arr] (b20)--(b21);
\end{tikzpicture}
\end{center}

Ruta corta (\textbf{7 bloques}) para quien quiere Go + infraestructura
sin pasar por C/Rust. Go es el lenguaje en el que están escritos Docker,
Kubernetes, Terraform, Prometheus, Grafana, etcd, containerd, Consul y
Vault --- la mayoría de herramientas cloud-native. Ideal para perfiles
SysAdmin que quieren leer el código fuente de sus herramientas, escribir
operadores de K8s, plugins de Terraform, o exporters de Prometheus.

\subsection{Ruta completa (secuencial sugerida)}

\begin{center}
\begin{tikzpicture}[
  node distance=0.3cm,
  blk/.style={rectangle, rounded corners=3pt, draw, thick,
    minimum width=1.2cm, minimum height=0.6cm,
    font=\scriptsize\bfseries},
  arr/.style={-{Stealth[length=4pt]}, thick}
]
\node[blk, fill=tTrack] (b01) {B01};
\node[blk, fill=tTrack, right=of b01] (b02) {B02};
\node[blk, fill=cTrack, right=of b02] (b03) {B03};
\node[blk, fill=qTrack, right=of b03] (b04) {B04};
\node[blk, fill=rTrack, right=of b04] (b05) {B05};
\node[blk, fill=cTrack, right=of b05] (b06) {B06};
\node[blk, fill=rTrack, right=of b06] (b07) {B07};
\node[blk, fill=lTrack, right=of b07] (b08) {B08};
\node[blk, fill=lTrack, right=of b08] (b09) {B09};
\node[blk, fill=lTrack, right=of b09] (b10) {B10};
\node[blk, fill=lTrack, right=of b10] (b11) {B11};
\node[blk, fill=rTrack, below=0.8cm of b01] (b12) {B12};
\node[blk, fill=rTrack, right=of b12] (b13) {B13};
\node[blk, fill=tTrack, right=of b13] (b14) {B14};
\node[blk, fill=tTrack, right=of b14] (b15) {B15};
\node[blk, fill=qTrack, right=of b15] (b16) {B16};
\node[blk, fill=qTrack, right=of b16] (b17) {B17};
\node[blk, fill=dTrack, right=of b17] (b18) {B18};
\node[blk, fill=dTrack, right=of b18] (b19) {B19};
\node[blk, fill=dTrack, right=of b19] (b20) {B20};
\node[blk, fill=dTrack, right=of b20] (b21) {B21};
\node[blk, fill=goTrack, below=0.8cm of b12] (b22) {B22};

\foreach \a/\b in {b01/b02,b02/b03,b03/b04,b04/b05,b05/b06,
  b06/b07,b07/b08,b08/b09,b09/b10,b10/b11}
  \draw[arr] (\a)--(\b);
\draw[arr] (b11.south) -- ++(0,-0.25) -| (b12.north);
\foreach \a/\b in {b12/b13,b13/b14,b14/b15,b15/b16,b16/b17,
  b17/b18,b18/b19,b19/b20,b20/b21}
  \draw[arr] (\a)--(\b);
\draw[arr] (b21.south) -- ++(0,-0.25) -| (b22.north);
\end{tikzpicture}
\end{center}

B01 $\to$ B02 $\to$ B03 $\to$ B04 $\to$ B05 $\to$ B06 $\to$ B07
$\to$ B08 $\to$ B09 $\to$ B10 $\to$ B11 $\to$ B12 $\to$ B13 $\to$
B14 $\to$ B15 $\to$ B16 $\to$ B17 $\to$ B18 $\to$ B19 $\to$ B20
$\to$ B21 $\to$ B22

% ===================================================================
\section{Leyenda de pistas temáticas}
% ===================================================================

\begin{itemize}[leftmargin=2cm]
  \item[\colorbox{tTrack}{\makebox[1.5cm]{}}]
    \textbf{Troncal} --- bloques fundacionales o que integran ambos
    lenguajes (B01, B02, B14, B15).
  \item[\colorbox{cTrack}{\makebox[1.5cm]{}}]
    \textbf{Pista C} --- programación en C puro (B03, B06).
  \item[\colorbox{rTrack}{\makebox[1.5cm]{}}]
    \textbf{Pista Rust} --- programación en Rust (B05, B07, B12, B13).
  \item[\colorbox{goTrack}{\makebox[1.5cm]{}}]
    \textbf{Pista Go} --- lenguaje dominante del ecosistema SysAdmin y
    cloud-native (B22). Docker, K8s, Terraform, Prometheus están
    escritos en Go.
  \item[\colorbox{lTrack}{\makebox[1.5cm]{}}]
    \textbf{Linux / Infraestructura} --- administración y servicios
    (B08, B09, B10, B11).
  \item[\colorbox{qTrack}{\makebox[1.5cm]{}}]
    \textbf{Herramientas / Calidad} --- build systems, patrones,
    testing (B04, B16, B17).
  \item[\colorbox{dTrack}{\makebox[1.5cm]{}}]
    \textbf{DevOps / Cloud} --- scripting, bases de datos, cloud,
    Kubernetes y observabilidad (B18, B19, B20, B21).
\end{itemize}

\end{document}
